\section{Supervision, Proactivité et Conclusions}

%------------------------------------------------------------
% Slide 16 : Approche de la Supervision et Volumétrie
%------------------------------------------------------------
\begin{frame}{Supervision Locale de l'Espace Disque}
    Pour préserver les ressources des postes, une supervision native et légère a été privilégiée.

    \vspace{0.2cm}
    \begin{columns}[T]
        \begin{column}{0.45\textwidth}
            \textbf{Analyse via \texttt{df -h} :}
            \begin{itemize}
                \item Surveillance du serveur de sauvegarde (\texttt{VLAN 60}).
                \item Utilisation des volumes logiques via LVM.
                \item \textbf{État actuel :} 4.6 Go utilisés (37\,\%), 8.1 Go disponibles.
                \item Point de contrôle clé pour éviter la saturation des données scientifiques.
            \end{itemize}
        \end{column}
        
        \begin{column}{0.55\textwidth}
            \begin{figure}
                \centering
                % Intégration de votre capture file_serveur_size.jpeg
                \includegraphics[width=0.98\textwidth]{figures/file_serveur_size.jpeg}
                \caption{Suivi de l'occupation des blocs du système de fichiers}
            \end{figure}
        \end{column}
    \end{columns}
\end{frame}

%------------------------------------------------------------
% Slide 17 : Innovation - ChatOps via Telegram Bot
%------------------------------------------------------------
% \begin{frame}{Alerte Proactive : Automatisation \& ChatOps}
%     Pour passer d'une maintenance réactive à un monitoring proactif sans alourdir le système.

%     \vspace{0.3cm}
%     \begin{block}{Le Script de Supervision Intelligent (Node.js)}
%         \begin{itemize}
%             \item Script local léger qui analyse les métriques de stockage en arrière-plan.
%             \item \textbf{Déclenchement immédiat} dès qu'un changement ou un seuil critique est détecté.
%             \item Connexion sécurisée via l'API Web et notification instantanée sur le canal \textbf{Telegram} de l'administrateur.
%         \end{itemize}
%     \end{block}

%     \begin{exampleblock}{Avantage Ingénieur}
%         L'administrateur réseau est alerté en temps réel sur son smartphone de la réussite ou de l'anomalie d'une réplication de sauvegarde.
%     \end{exampleblock}
% \end{frame}

%------------------------------------------------------------
% Slide 18 : Retours d'Expérience & Défis Surmontés
%------------------------------------------------------------
\begin{frame}{Retours d'Expérience et Défis Surmontés}
    La réalisation de cette infrastructure a imposé des arbitrages techniques forts face aux réalités logistiques.

    \vspace{0.2cm}
    \begin{itemize}
        \item \textbf{Gestion du temps :} Le calendrier global ne nous ayant pas été favorable, l'équipe a dû se mobiliser en mode Sprint Agile pour prioriser les livrables vitaux.
        \item \textbf{Limitations matérielles :} La charge CPU/RAM générée par l'émulation complète sous GNS3 a ralenti nos postes de travail.
        \item \textbf{Solution et posture MVP :} Choix d'une architecture rationnelle centrée sur un serveur d'infrastructure robuste et un nœud de sauvegarde étanche.
    \end{itemize}
\end{frame}

%------------------------------------------------------------
% Slide 19 : Perspectives d'Évolution
%------------------------------------------------------------
\begin{frame}{Perspectives d'Évolution de l'Infrastructure}
    Pour accompagner la croissance future de l'institut \textit{DataLab Research} :

    \vspace{0.3cm}
    \begin{columns}[T]
        \begin{column}{0.31\textwidth}
            \begin{block}{Ansible}
                \small
                Industrialisation via l'**Infrastructure as Code** (Playbooks) pour standardiser les configurations.
            \end{block}
        \end{column}

        \begin{column}{0.31\textwidth}
            \begin{block}{Zabbix}
                \small
                Centralisation des métriques sur une interface graphique unifiée avec alertes visuelles.
            \end{block}
        \end{column}

        \begin{column}{0.31\textwidth}
            \begin{block}{Docker}
                \small
                Conteneurisation des services pour économiser jusqu'à 80\,\% de RAM sur la maquette.
            \end{block}
        \end{column}
    \end{columns}
\end{frame}

%------------------------------------------------------------
% Slide 20 : Conclusion & Page Finale
%------------------------------------------------------------
\begin{frame}[plain]
    \centering
    \vspace{2cm}
    {\huge \color{blue!70!black} \textbf{Merci pour votre attention !}}
    
    \vspace{1cm}
    {\large Place aux questions du jury.}
    
    \vspace{1.5cm}
    \begin{block}{Code Source et Scripts du Projet}
        \centering
        \tiny \url{https://github.com/Jovialngandu/datalab_infra_m1}
    \end{block}
\end{frame}